% Chapter 10 — An ML Implementation of Simple Types
% Source: ../source/split/chapters/ch10-ml-simple-types_p113-p116.pdf
% Original print pages: 113--116
% Status: 已确认

\chapter{简单类型的 ML 实现}
\label{chap:10}

把 \(\lambda_{\to}\) 具体实现为 ML 程序，所遵循的思路与\chapref{7}
实现无类型 lambda 演算时相同。主要增加的是函数 \texttt{typeof}，它在给定
上下文中计算给定项的类型。不过在介绍它之前，还需要一些操纵上下文的底层
机制。

\translatornote{原书的 OCaml 程序完整保留在正文中，每段代码之后直接给出
译者增补的 Rust 对照实现。Rust 版本不追求逐行照译，而是保持相同的对象语言
语法、类型检查规则和求值行为。}

\section{上下文}
\label{sec:10.1}

回想 \taplsecref{7.1}{sec:7.1}，上下文是一个列表；列表中的每一项都是一个
二元组，包含变量名及其绑定信息：
\begin{taplocamlcode}
type context = (string * binding) list
\end{taplocamlcode}
\taplrustcounterpart{ch10-context}

在\chapref{7}中，上下文只服务于两项工作：解析时把项从有名形式转换为
无名形式，打印时再把无名形式还原为有名形式。为此，我们只需记录变量名；
\texttt{binding} 类型被定义成一个只有单个构造子的平凡数据类型，不携带
任何信息：
\begin{taplocamlcode}
type binding = NameBind
\end{taplocamlcode}
\taplrustcounterpart{ch10-binding-name}

实现类型检查器时\footnote{这里描述的实现对应带布尔值（\cref{fig:8-1}）的
简单类型 lambda 演算（\cref{fig:9-1}）。本章代码可在作者网站仓库的
\texttt{simplebool} 实现中找到。}，需要用上下文携带关于变量的类型假设。因此，我们为
\texttt{binding} 类型加入一个名为 \texttt{VarBind} 的新构造子：
\begin{taplocamlcode}
type binding =
    NameBind
  | VarBind of ty
\end{taplocamlcode}
\taplrustcounterpart{ch10-binding-var}

每个 \texttt{VarBind} 构造子都携带相应变量的类型假设。旧的
\texttt{NameBind} 仍然保留，供不关心类型假设的打印和解析函数使用。
（另一种实现策略是定义两套完全不同的上下文类型：一套用于解析和打印，
另一套用于类型检查。）

\texttt{typeof} 使用 \texttt{addbinding} 函数，把新的变量绑定
\texttt{(x,bind)} 加入上下文 \texttt{ctx}。上下文由列表表示，所以
\texttt{addbinding} 实质上就是列表的 cons 操作：
\begin{taplocamlcode}
let addbinding ctx x bind = (x,bind)::ctx
\end{taplocamlcode}
\taplrustcounterpart{ch10-add-binding}

反过来，函数 \texttt{getTypeFromContext} 从上下文 \texttt{ctx} 中提取
与特定变量 \texttt{i} 关联的类型假设。若 \texttt{i} 越界，则使用文件
位置信息 \texttt{fi} 打印错误消息：
\begin{taplocamlcode}
let getTypeFromContext fi ctx i =
  match getbinding fi ctx i with
      VarBind(tyT) -> tyT
    | _ -> error fi
      ("getTypeFromContext: Wrong kind of binding for variable "
      ^ (index2name fi ctx i))
\end{taplocamlcode}
\taplrustcounterpart{ch10-get-type-from-context}

这个模式匹配还提供了一项内部一致性检查：正常情况下，调用
\texttt{getTypeFromContext} 时，上下文中的第 \texttt{i} 个绑定应当确实
是 \texttt{VarBind}。不过，后续章节会加入其他形式的绑定，特别是类型变量
的绑定。届时，如果传入的索引指向的不是 \texttt{VarBind}，
\texttt{getTypeFromContext} 就无法从中取得项变量的类型，于是会调用底层的
\texttt{error} 函数报告这种内部不一致；传入的 \texttt{info} 用于指出错误
发生的文件位置。
\begin{taplocamlcode}
val error : info -> string -> 'a
\end{taplocamlcode}
\taplrustcounterpart{ch10-error}

\texttt{error} 的结果类型是\tapltranslateditalic{类型变量} \texttt{'a}，可以实例化为
任意 ML 类型。这是合理的，因为该函数无论如何也不会返回：它打印消息后便
终止程序。这里必须假定 \texttt{error} 的结果是 \texttt{ty}，因为模式
匹配的另一分支返回的正是 \texttt{ty}。

\translatornote{Rust 对照实现把 \texttt{error} 的返回类型写成
\texttt{!}，表示该函数永远不会正常返回。因此，对 \texttt{error} 的调用
也可以出现在需要任何其他结果类型的位置；在这里，\texttt{!} 与 OCaml 中
可以实例化为任意类型的 \texttt{'a} 起到相同作用。}

注意，类型假设是按索引查找的，因为项在内部采用无名形式表示，变量表示为
数字索引。函数 \texttt{getbinding} 只是查找给定上下文中的第
\texttt{i} 个绑定：
\begin{taplocamlcode}
val getbinding : info -> context -> int -> binding
\end{taplocamlcode}
\taplrustcounterpart{ch10-get-binding}
其定义可在本书网站的 \texttt{simplebool} 实现中找到。

\section{项与类型}
\label{sec:10.2}

类型语法直接从\cref{fig:8-1} 与\cref{fig:9-1} 的抽象语法转写为
ML 数据类型：
\begin{taplocamlcode}
type ty =
    TyBool
  | TyArr of ty * ty
\end{taplocamlcode}
\taplrustcounterpart{ch10-type}

项的表示与无类型 lambda 演算所用的形式相同
（见 \taplsecref{7.1}{sec:7.1}），只是在
\texttt{TmAbs} 分支中加入类型标注：
\begin{taplocamlcode}
type term =
    TmTrue of info
  | TmFalse of info
  | TmIf of info * term * term * term
  | TmVar of info * int * int
  | TmAbs of info * string * ty * term
  | TmApp of info * term * term
\end{taplocamlcode}
\taplrustcounterpart{ch10-term}

\section{类型检查}
\label{sec:10.3}

类型检查函数 \texttt{typeof} 可以看作 \(\lambda_{\to}\) 类型规则
（\cref{fig:8-1} 与\cref{fig:9-1}）的直接翻译；更准确地说，它是
\taplnamedref{lem:9.3.1}{反演引理}的转写。后一种说法更准确，是因为反演引理针对每一种
句法形式，精确说明这种形式的项成为良类型项必须满足什么条件。类型规则只
说明特定形式的项在特定条件下是良类型的；单看一条类型规则，永远不能断定
某个项不是良类型的，因为总可能有另一条规则能为这个项赋型。（目前这看起来
也许只是没有实质区别的措辞差异，因为反演引理直接来自类型规则。但在后面的
系统中，证明反演引理会比在 \(\lambda_{\to}\) 中费力得多，这项区别便会
变得重要。）

\begin{taplocamlcode}
let rec typeof ctx t =
  match t with
    TmTrue(fi) ->
      TyBool
  | TmFalse(fi) ->
      TyBool
  | TmIf(fi,t1,t2,t3) ->
      if (=) (typeof ctx t1) TyBool then
        let tyT2 = typeof ctx t2 in
        if (=) tyT2 (typeof ctx t3) then tyT2
        else error fi "arms of conditional have different types"
      else error fi "guard of conditional not a boolean"
  | TmVar(fi,i,_) ->
      getTypeFromContext fi ctx i
  | TmAbs(fi,x,tyT1,t2) ->
      let ctx' = addbinding ctx x (VarBind(tyT1)) in
      let tyT2 = typeof ctx' t2 in
      TyArr(tyT1, tyT2)
  | TmApp(fi,t1,t2) ->
      let tyT1 = typeof ctx t1 in
      let tyT2 = typeof ctx t2 in
      (match tyT1 with
         TyArr(tyT11,tyT12) ->
           if (=) tyT2 tyT11 then tyT12
           else error fi "parameter type mismatch"
       | _ -> error fi "arrow type expected")
\end{taplocamlcode}
\taplrustcounterpart{ch10-type-of}

这里有两个 OCaml 细节值得一提。第一，等号运算符 \texttt{=} 被写在括号中，
因为这里把它放在前缀位置使用，而不是通常的中缀位置。这样更方便与后续版本的
\texttt{typeof} 比较；到那时，比较类型需要比简单相等更精细的操作。

第二，等号运算符计算复合值的结构相等，而不是指针相等。也就是说：
\enlargethispage{3\baselineskip}
\begin{taplocamlcode}
let t = TmApp(t1,t2) in
let t' = TmApp(t1,t2) in
(=) t t'
\end{taplocamlcode}
虽然绑定到 \texttt{t} 与 \texttt{t'} 的两个 \texttt{TmApp} 实例是在不同
时刻分别分配的，内存地址也不同，上述比较仍会得到 \texttt{true}。
\taplrustexplanation{Rust 不需要为这个例子另写相等性函数。\texttt{Term}
已经派生 \texttt{PartialEq} 与 \texttt{Eq}；因此，对两个分别构造但内容
相同的项直接使用 \texttt{t == t'}，就会递归比较其结构和内容。这与上面的
OCaml 表达式 \texttt{(=) t t'} 作用相同。}
